\documentclass[aspectratio=169,11pt]{beamer}

% ============================================================
% PACOTES
% ============================================================
\usepackage[utf8]{inputenc}
\usepackage[T1]{fontenc}
\usepackage[brazil]{babel}
\usepackage{amsmath,amssymb}
\usepackage{xcolor}

% ============================================================
% TEMA
% ============================================================
\usetheme{Madrid}
\usecolortheme{whale}
\setbeamertemplate{navigation symbols}{}
\setbeamertemplate{footline}[frame number]

\definecolor{darkblue}{RGB}{0,51,102}
\setbeamercolor{title}{fg=white,bg=darkblue}
\setbeamercolor{frametitle}{fg=white,bg=darkblue}

\title[Guia --- Decisões de LLMs]{Guia do Apresentador\\Explicando Decisões de LLMs via Otimização Inversa}
\author{George Gomes}
\date{2026}

\begin{document}

% ============================================================
% SLIDE 1
% ============================================================
\begin{frame}{Slide 1 --- Título e Identificação}
\small
Bom dia a todos, meu nome é George Gomes e esta é a defesa da minha dissertação de mestrado, intitulada ``Explicando Decisões de LLMs via Otimização Inversa''.

Este trabalho investiga se modelos de linguagem de grande escala --- os LLMs --- possuem um processo de tomada de decisão que pode ser capturado matematicamente. A ideia central é usar técnicas de otimização inversa para extrair uma métrica que explique como o LLM decide.

Vou estruturar a apresentação em 9 partes: motivação, design experimental, os resultados das 4 fases do experimento, análises de viés e sensibilidade, validações adicionais, análise estatística, e conclusões.
\end{frame}

% ============================================================
% SLIDE 2
% ============================================================
\begin{frame}{Slide 2 --- Motivação}
\small
O ponto de partida deste trabalho é uma observação simples mas importante: LLMs tomam decisões o tempo todo --- classificam textos, respondem perguntas, fazem julgamentos --- mas não explicam qual critério estão usando internamente.

A pergunta central que nos guia é: ``LLMs possuem um processo de decisão estável e aprendível?'' Se a resposta for sim, então podemos extrair uma métrica matemática W que funcione como uma ``explicação'' das decisões do modelo. Isso tem implicações diretas para explainability, para construção de confiança em sistemas de IA, e para auditoria de decisões automatizadas.

A relevância prática é clara: se eu consigo descrever matematicamente como o LLM decide, posso prever, auditar e eventualmente corrigir essas decisões.
\end{frame}

% ============================================================
% SLIDE 3
% ============================================================
\begin{frame}{Slide 3 --- Hipóteses}
\small
Formulamos cinco hipóteses que guiam todo o trabalho experimental.

H1 afirma que o LLM mantém o mesmo critério decisório quando enfrenta problemas com geometrias diferentes. H2 diz que fornecer exemplos rotulados no prompt aumenta a concordância entre o LLM e a métrica aprendida. H3 inverte o papel: pergunta se o LLM consegue aprender o critério de um perito externo via few-shot learning. H4 é particularmente interessante --- postulamos que exemplos próximos à fronteira de decisão seriam mais informativos que exemplos fáceis. E H5 testa se tudo isso é reproduzível entre diferentes execuções com sementes aleatórias distintas.

Já adianto que nem todas foram confirmadas, e a refutação de H4 é um dos achados mais relevantes do trabalho.
\end{frame}

% ============================================================
% SLIDE 4
% ============================================================
\begin{frame}{Slide 4 --- Abordagem (Otimização Inversa)}
\small
A abordagem é conceitualmente elegante: dado que o LLM tomou decisões observáveis, tentamos descobrir qual função objetivo tornaria essas decisões ótimas. Isso é exatamente o que a otimização inversa faz.

A métrica que buscamos é uma distância de Mahalanobis diagonal, onde cada peso $w_i$ indica a importância relativa de cada feature. Escolhemos a forma diagonal deliberadamente: com $O(d)$ parâmetros em vez de $O(d^2)$, ganhamos interpretabilidade e tratabilidade computacional. Cada $w_i$ tem significado direto.

A classificação é do tipo nearest-centroid: cada ponto é atribuído ao centroide mais próximo segundo a métrica W. Isso é simples, mas suficiente para capturar critérios anisotrópicos --- ou seja, critérios que pesam as dimensões de forma diferente.
\end{frame}

% ============================================================
% SLIDE 5
% ============================================================
\begin{frame}{Slide 5 --- Dados Sintéticos 2D}
\small
Usamos dados sintéticos bidimensionais gerados com a função make\_blobs do scikit-learn. São 4 problemas distintos com distribuições Gaussianas.

A escolha por dados 2D é proposital: permite visualização direta das fronteiras de decisão, dos pontos classificados corretamente e incorretamente, e da geometria aprendida. Cada problema tem centroides e espalhamento diferentes, criando geometrias distintas para testar a generalização da métrica.

No gráfico à direita vocês podem ver os scatter plots dos problemas lado a lado. As cores representam as classes verdadeiras do ground truth.
\end{frame}

% ============================================================
% SLIDE 6
% ============================================================
\begin{frame}{Slide 6 --- Problema A}
\small
O Problema A é o problema de treino. Tem 150 pontos divididos em 2 classes. É neste problema que o LLM faz suas classificações em zero-shot --- sem nenhum exemplo --- e a partir dessas 150 decisões, nossos algoritmos de otimização inversa inferem a métrica W que melhor explica as escolhas do LLM.

Os parâmetros dos centroides e do cluster\_std foram escolhidos para criar classes com alguma sobreposição, mas não excessiva. Isso é importante: se os clusters fossem perfeitamente separados, qualquer métrica funcionaria; a sobreposição é o que nos permite distinguir métricas diferentes.
\end{frame}

% ============================================================
% SLIDE 7
% ============================================================
\begin{frame}{Slide 7 --- Problema B}
\small
O Problema B tem 100 pontos e centroides deslocados verticalmente em relação ao Problema A. A geometria é diferente, mas a tarefa é a mesma: classificação binária em 2D.

Este é o primeiro teste de consistência: aplicamos a métrica aprendida no Problema A e verificamos se o LLM, classificando pontos do Problema B, concorda com as predições dessa métrica transferida. Se concordar, temos evidência de que o critério decisório do LLM é estável e transferível entre problemas.
\end{frame}

% ============================================================
% SLIDE 8
% ============================================================
\begin{frame}{Slide 8 --- Problema C}
\small
O Problema C é semelhante ao B, mas com distorção geométrica mais severa --- 100 pontos com uma configuração mais desafiadora. Funciona como um segundo nível de teste: se a métrica generaliza para o Problema B mas não para o C, sabemos que a consistência tem limites e depende de quão diferente é a geometria.

A combinação B+C nos dá uma visão gradual: consistência moderada a difícil.
\end{frame}

% ============================================================
% SLIDE 9
% ============================================================
\begin{frame}{Slide 9 --- Problema D}
\small
O Problema D é fundamentalmente diferente dos anteriores. Aqui, temos um expert externo com uma métrica W conhecida --- especificamente, W = [0.3, 1.5], o que significa que a segunda dimensão, $x_2$, é 5 vezes mais importante que $x_1$ na classificação.

O expert rotula os 150 pontos do Problema D usando essa métrica. O papel do LLM aqui é invertido: em vez de ser o ``sujeito'' cujas decisões analisamos, o LLM é o ``aprendiz'' que deve aprender o critério do expert a partir de exemplos. No gráfico à direita vocês podem ver o Problema D com a fronteira de decisão do expert claramente desenhada.
\end{frame}

% ============================================================
% SLIDE 10
% ============================================================
\begin{frame}{Slide 10 --- Configurações de Execução}
\small
Um dos pilares do trabalho é o design de variabilidade. Precisamos garantir que os resultados não são artefatos de uma configuração específica.

Usamos 3 sementes aleatórias diferentes --- 42, 123 e 7 --- que controlam a geração dos dados. Para cada configuração, fazemos 3 repetições. A temperatura do modelo está fixa em 0.0 para minimizar a estocasticidade --- queremos isolar o critério decisório, não o ruído da amostragem do modelo. Usamos o gpt-4o-mini via API da OpenAI.

Para detectar viés semântico, testamos 4 pares de nomes de classe: A/B que é neutro, 0/1 que é numérico, Positivo/Negativo e Azul/Vermelho que são semânticos. A ideia é verificar se o nome da classe influencia a decisão do LLM --- e, como veremos, influencia significativamente.
\end{frame}

% ============================================================
% SLIDE 11
% ============================================================
\begin{frame}{Slide 11 --- Perceptron Estruturado}
\small
O primeiro algoritmo é um Perceptron Estruturado com relaxação de margem. A ideia é buscar iterativamente os pesos W que maximizem a separação entre as classes segundo a métrica de Mahalanobis diagonal.

O procedimento usa busca binária sobre o parâmetro gamma, que é a margem mínima desejada. Para cada valor de gamma, treinamos o perceptron até convergência: quando a métrica erra --- ou seja, atribui um ponto à classe errada segundo o LLM --- atualizamos os pesos na direção que corrige o erro. Os pesos são projetados para não-negatividade após cada atualização.

O resultado é o W com o maior gamma que ainda permite convergência. Quanto maior o gamma, mais robusta é a separação --- mas se as decisões do LLM forem inconsistentes, gamma será pequeno ou zero.
\end{frame}

% ============================================================
% SLIDE 12
% ============================================================
\begin{frame}{Slide 12 --- Mínimos Quadrados (NNLS)}
\small
O segundo algoritmo aborda o problema de forma diferente: formula as restrições de classificação correta como um sistema de desigualdades lineares e resolve via Non-Negative Least Squares.

Para cada ponto classificado como classe $k$ pelo LLM, escrevemos a restrição de que a distância à classe correta deve ser menor que a distância a qualquer outra classe. Isso gera um sistema $Aw \leq b$ que resolvemos minimizando o resíduo quadrático com a restrição $w \geq 0$.

A grande vantagem é que a solução não é iterativa --- é computada diretamente. E o fato de termos dois algoritmos independentes nos dá uma validação cruzada: se ambos convergem para um W similar, o resultado é mais confiável.
\end{frame}

% ============================================================
% SLIDE 13
% ============================================================
\begin{frame}{Slide 13 --- Programação Linear (LP)}
\small
O terceiro algoritmo é uma formulação de programação linear que maximiza a margem mínima sobre todas as observações. A ideia é encontrar o W que, sujeito a classificar todos os pontos corretamente segundo a métrica, maximiza a menor margem individual.

A normalização de que a soma dos pesos é igual a 1 evita a solução trivial de aumentar todos os pesos indefinidamente. A solução é via o solver HiGHS do scipy.

Quando gamma ótimo é zero, significa que não existe nenhuma métrica diagonal que separe perfeitamente os dados segundo as decisões do LLM. Na prática, isso aconteceu em vários cenários da nossa execução, o que indica que o LLM não é perfeitamente consistente com nenhuma métrica de Mahalanobis diagonal.
\end{frame}

% ============================================================
% SLIDE 14
% ============================================================
\begin{frame}{Slide 14 --- Fase A Descrição}
\small
Vamos agora aos resultados. Na Fase A, o LLM classifica 150 pontos do Problema A em zero-shot. O prompt é extremamente simples: apresenta as coordenadas do ponto e pede que o modelo classifique como A ou B.

A partir dessas 150 decisões, cada um dos 3 algoritmos infere uma métrica W. É importante frisar que chamamos isso de ``W estimado'' ou ``W inferido'', não ``W aprendido pelo LLM'' --- o aprendizado é nosso, via otimização inversa, sobre as decisões observadas do LLM.

O W resultante representa o critério decisório implícito que melhor explica as classificações feitas pelo modelo.
\end{frame}

% ============================================================
% SLIDE 15
% ============================================================
\begin{frame}{Slide 15 --- Fase A --- W e Fidelidade}
\small
Aqui temos os resultados concretos. Para seed=42 com classes A/B, o Perceptron estimou W unitário como [1.0, 0.0], ou seja, só a primeira dimensão importa. O NNLS estimou [0.20, 0.80], priorizando a segunda dimensão. O LP ficou em [0.71, 0.71], próximo da euclidiana.

A fidelidade --- ou seja, a porcentagem de decisões do LLM que a métrica reproduz no próprio Problema A --- ficou entre 72.7\% e 80.7\%. Isso não é perfeito, e sinaliza que o LLM não é totalmente consistente com nenhuma métrica diagonal simples.

Um detalhe importante: a acurácia do LLM versus o ground truth é apenas 30.7\%. Isso parece ruim, mas lembrem que estamos medindo se a métrica reproduz o LLM, não se o LLM acerta. Um LLM pode ter um critério estável que simplesmente discorda do ground truth.
\end{frame}

% ============================================================
% SLIDE 16
% ============================================================
\begin{frame}{Slide 16 --- Distribuição de W entre Seeds}
\small
Este é um resultado preocupante para H5. Quando olhamos a distribuição de W estimado entre as 3 seeds, vemos instabilidade significativa. Na seed 42, o Perceptron estimou W apontando na direção de $x_1$; nas seeds 123 e 7, apontou na direção de $x_2$ --- praticamente ortogonais.

A similaridade cosseno média entre seeds é apenas 0.53, e o coeficiente de variação da razão $w_1/w_2$ é de 258.9\% --- extremamente instável. Isso sugere que o critério decisório zero-shot do LLM pode não ser estável o suficiente para ser capturado de forma confiável.

No entanto, isso pode refletir que o LLM, sem orientação, simplesmente não tem um critério forte e usa heurísticas que variam com a disposição espacial dos dados.
\end{frame}

% ============================================================
% SLIDE 17
% ============================================================
\begin{frame}{Slide 17 --- Comparação de Algoritmos}
\small
Comparando os 3 algoritmos, vemos que NNLS e LP são mais concordantes entre si, com similaridade cosseno de 0.87. O Perceptron diverge mais em certas seeds, especialmente quando os dados não são linearmente separáveis sob nenhuma métrica diagonal --- nesses casos, ele atinge soluções de ``melhor esforço'' que podem ser degeneradas.

A fidelidade média confirma que o NNLS é ligeiramente superior: 82.4\% contra 76\% do Perceptron e 75.2\% do LP. O LP frequentemente retorna a solução isotrópica [0.5, 0.5] quando o problema é infeasible, o que é uma forma de ``fallback'' para a métrica euclidiana.

A divergência entre algoritmos não invalida a abordagem --- pelo contrário, mostra que diferentes formulações capturam aspectos diferentes da inconsistência do LLM.
\end{frame}

% ============================================================
% SLIDE 18
% ============================================================
\begin{frame}{Slide 18 --- Fases B/C Descrição}
\small
As Fases B e C testam a pergunta central de H1: a métrica aprendida na Fase A generaliza para problemas com geometria diferente?

O procedimento é direto: pegamos o W estimado no Problema A, calculamos centróides locais nos Problemas B e C, e comparamos as predições da métrica com as classificações do LLM nesses novos problemas.

Usamos três métricas de concordância: a consistência bruta em porcentagem, o Cohen's Kappa que ajusta pelo acaso, e o F1-Score. O Kappa é particularmente importante porque leva em conta a distribuição das classes --- um classificador que sempre diz ``A'' pode ter alta consistência bruta com outro que também tende a dizer ``A'', mas o Kappa revelará que a concordância é espúria.
\end{frame}

% ============================================================
% SLIDE 19
% ============================================================
\begin{frame}{Slide 19 --- Fases B/C Resultados}
\small
Os resultados aqui são muito reveladores. Em zero-shot, a consistência é fraca: 64\% no Problema B e 63\% no C, com Kappa de apenas 0.14. Isso significa que, sem exemplos, o LLM e a métrica concordam pouco mais do que o acaso.

Mas o salto com few-shot é dramático: com apenas 5 exemplos rotulados pela métrica, a consistência sobe para 91.7\% e 92.7\%, com Kappa acima de 0.80 --- concordância substancial segundo os critérios de Landis e Koch. Com 10 exemplos, chegamos a 96.3\% e Kappa de 0.93.

A saturação ocorre por volta de 10-20 exemplos; adicionar mais não melhora significativamente. Isso mostra que poucos exemplos bem escolhidos são suficientes para ``calibrar'' o LLM com a métrica.
\end{frame}

% ============================================================
% SLIDE 20
% ============================================================
\begin{frame}{Slide 20 --- Fases B/C Discussão}
\small
Vamos à interpretação. H1, sobre consistência, é parcialmente confirmada: o LLM mantém consistência elevada com a métrica quando recebe exemplos, mas não em zero-shot. Isso sugere que o LLM tem a \textit{capacidade} de seguir um critério consistente, mas precisa de orientação para ativá-lo.

H2 é claramente confirmada: exemplos rotulados pela métrica amplificam enormemente a concordância. O salto de zero-shot para 5-shot é de aproximadamente 30 pontos percentuais --- um efeito muito grande.

Uma limitação importante é a fidelidade baixa na Fase A: se a métrica já não reproduz perfeitamente o LLM no Problema A, a base para transferência é frágil. Isso levanta uma questão sobre circularidade no few-shot: os exemplos são rotulados pela própria métrica, então estamos essencialmente ``ensinando'' o LLM a concordar conosco.
\end{frame}

% ============================================================
% SLIDE 21
% ============================================================
\begin{frame}{Slide 21 --- Viés de Ordem das Classes}
\small
Uma análise de sensibilidade crucial é o viés de ordem das classes. Testamos se trocar ``A ou B'' por ``B ou A'' no prompt --- ou seja, simplesmente inverter a ordem em que as classes são apresentadas --- muda as decisões do LLM.

Os resultados mostram que sim, existe viés posicional. O exemplo mais dramático é o par A/B versus B/A: a consistência no Problema B salta de 79.5\% para 94.9\% quando invertemos a ordem. Isso é uma diferença de 15.4 pontos percentuais, causada apenas pela ordem das palavras no prompt.

Esse achado é relevante porque mostra que a métrica W inferida pode ser parcialmente um artefato da formulação do prompt, não apenas do critério decisório intrínseco do LLM. É um confundidor que precisa ser reportado e controlado.
\end{frame}

% ============================================================
% SLIDE 22
% ============================================================
\begin{frame}{Slide 22 --- Nomes de Features}
\small
Outro teste de sensibilidade: o que acontece quando mudamos os nomes das coordenadas de ``x1, x2'' para ``altura, peso'' ou ``feature\_1, feature\_2''?

Nomes como ``altura'' e ``peso'' são semanticamente carregados --- o LLM pode ter conhecimento prévio sobre a relação entre essas variáveis, o que influenciaria seus pesos implícitos. Por exemplo, pode ``saber'' que peso é mais discriminativo para certas classificações.

Os resultados confirmam que o LLM é sensível aos nomes das features. O W estimado muda dependendo dos nomes usados, o que é mais uma evidência de que o critério decisório do LLM é influenciado pelo contexto semântico do prompt, não apenas pelos dados numéricos.
\end{frame}

% ============================================================
% SLIDE 23
% ============================================================
\begin{frame}{Slide 23 --- Variantes de Prompt (Design)}
\small
Para quantificar mais formalmente a sensibilidade ao prompt, testamos 4 variantes de template. A variante default é o nosso template baseline. A variante geometric adiciona contexto espacial explícito, como ``considere a distância entre os pontos''. A variante chain-of-thought pede que o LLM raciocine passo a passo antes de responder. E a variante tabular apresenta as coordenadas como uma tabela markdown em vez de texto corrido.

Cada variante executa as Fases A, B e C completas em todas as seeds, o que nos dá uma comparação controlada. A análise mais importante é o scatter de W aprendidos: se os pontos se agrupam por variante e não por seed, o prompt é um confundidor forte.
\end{frame}

% ============================================================
% SLIDE 24
% ============================================================
\begin{frame}{Slide 24 --- Variantes de Prompt (Resultados)}
\small
Os resultados mostram sensibilidade moderada ao template do prompt. O W estimado varia entre variantes, com a variante chain-of-thought e a geometric tendendo a alterar a distribuição de decisões do LLM em relação ao baseline.

A implicação é séria: se a métrica muda com o prompt, até que ponto estamos medindo algo intrínseco ao LLM versus algo que depende de como formulamos a pergunta? Isso é análogo a um instrumento de medição que produz resultados diferentes dependendo de como você segura ele.

No entanto, a consistência few-shot se mantém alta em todas as variantes, o que sugere que a capacidade do LLM de seguir um critério com orientação é mais robusta que o critério zero-shot em si.
\end{frame}

% ============================================================
% SLIDE 25
% ============================================================
\begin{frame}{Slide 25 --- Fase D Descrição}
\small
Agora entramos na Fase D, que é conceitualmente diferente de tudo anterior. Aqui o papel se inverte: o LLM não é o sujeito de análise, mas o aprendiz.

Temos um expert externo com uma métrica W conhecida --- $[0.3, 1.5]$, que dá peso 5 vezes maior à segunda dimensão. O expert rotula os pontos do Problema D e fornece exemplos ao LLM via few-shot. A pergunta é simples: com quantos exemplos o LLM consegue reproduzir as classificações do expert? E importa como esses exemplos são escolhidos?

No gráfico do Problema D com a fronteira do expert, vocês podem ver que a fronteira é inclinada --- refletindo a anisotropia do W. Pontos longe da fronteira são ``fáceis'' e pontos próximos são ``difíceis''.
\end{frame}

% ============================================================
% SLIDE 26
% ============================================================
\begin{frame}{Slide 26 --- Estratégias de Seleção}
\small
Testamos 4 estratégias para escolher quais exemplos fornecer ao LLM. A estratégia Easy seleciona pontos longe da fronteira de decisão --- são os pontos que o expert classifica com alta confiança. A estratégia Hard seleciona pontos próximos à fronteira --- são os ambíguos, que em princípio carregam mais informação sobre onde está a fronteira. A estratégia Mixed combina 50\% de cada. E Random é o baseline sem critério.

No gráfico à direita, vocês podem ver a distribuição espacial: os pontos Easy estão nas extremidades dos clusters, e os Hard estão aglomerados na região de transição entre as classes.

A intuição de H4 era que pontos Hard, por estarem na fronteira, informariam melhor o LLM sobre a geometria da decisão. Vamos ver se essa intuição se confirma.
\end{frame}

% ============================================================
% SLIDE 27
% ============================================================
\begin{frame}{Slide 27 --- Curvas de Aprendizado}
\small
Aqui temos o resultado mais importante da Fase D: as curvas de aprendizado. A tabela mostra a acurácia do LLM versus o expert para cada número de exemplos e cada estratégia.

O ponto de partida é 55.5\% em zero-shot --- pouco melhor que o aleatório. Com 5 exemplos easy, já sobe para 78.6\%. Com 40 exemplos hard, chega a 95.8\%.

H3 está claramente confirmada: o LLM consegue aprender o critério do expert, e melhora consistentemente com mais exemplos. A melhoria total é de mais de 40 pontos percentuais.

Mas observem o comportamento da estratégia Hard: com 5 exemplos, fica em 52.5\% --- pior que zero-shot! Só melhora significativamente a partir de 20 exemplos. Isso é muito diferente de Easy e Random, que já são bons com poucos exemplos.
\end{frame}

% ============================================================
% SLIDE 28
% ============================================================
\begin{frame}{Slide 28 --- Comparação de Estratégias}
\small
Agrupando os resultados por estratégia, a média de acurácia para todos os tamanhos de few-shot maiores que zero é: Random com 90.4\%, Mixed com 88.1\%, Easy com 81.4\%, e Hard com 75.3\%.

H4 é refutada. Exemplos fáceis superam exemplos difíceis consistentemente. E mais: a estratégia aleatória é a melhor! Isso parece contraintuitivo, mas faz sentido quando pensamos em como LLMs processam exemplos in-context.

A interpretação é que exemplos fáceis funcionam como ``âncoras'' que definem claramente as regiões de cada classe. O LLM consegue interpolar a partir dessas âncoras para classificar pontos intermediários. Exemplos difíceis, por estarem na zona de ambiguidade, não fornecem padrões claros de nenhuma das classes --- confundem mais do que ajudam, especialmente com poucos exemplos.

Este é um achado que consideramos relevante para a literatura de in-context learning.
\end{frame}

% ============================================================
% SLIDE 29
% ============================================================
\begin{frame}{Slide 29 --- Múltiplos Experts}
\small
Para verificar se os resultados dependem do expert específico, testamos 3 configurações: aniso\_x2 com W = [0.3, 1.5] que é o original, aniso\_x1 com W = [1.5, 0.3] que inverte a importância das dimensões, e euclidean com W = [1.0, 1.0] que dá pesos iguais.

Isso responde à pergunta: ``O LLM aprende igualmente bem critérios diferentes?''

Como validação, rodamos os 3 algoritmos sobre os rótulos do oráculo --- ou seja, sobre as decisões de um classificador que usa exatamente o W do expert. Os 3 algoritmos recuperam o W com similaridade cosseno de 0.99 a 1.00 e fidelidade de 95\% a 100\%. Isso confirma que os algoritmos estão corretos: quando o classificador é perfeitamente consistente, o W é recuperado com precisão.
\end{frame}

% ============================================================
% SLIDE 30
% ============================================================
\begin{frame}{Slide 30 --- Diluição}
\small
O experimento de diluição testa uma preocupação prática: se adicionarmos exemplos fáceis a um conjunto de exemplos difíceis, os fáceis ``diluem'' a informação dos difíceis?

O design fixa 4 exemplos hard e adiciona progressivamente 0, 2, 4, 10, 16 e 20 exemplos easy. Os resultados mostram que a acurácia melhora monotonicamente: de 56.8\% com 4 hard puros para 86.9\% com 4 hard + 16 easy. Há uma leve queda para 84.9\% com 20 easy, mas não é significativa.

A conclusão é que exemplos fáceis não diluem --- eles complementam. Isso reforça o achado de que exemplos fáceis têm valor informativo significativo para o LLM, mesmo na presença de exemplos difíceis.
\end{frame}

% ============================================================
% SLIDE 31
% ============================================================
\begin{frame}{Slide 31 --- Viés de Ordem dos Exemplos}
\small
Este é um teste específico da Fase D: a ordem em que os exemplos few-shot são apresentados afeta o resultado? Isso é relevante por causa do chamado ``recency bias'' --- a tendência de LLMs a dar mais peso aos últimos tokens processados.

Testamos 4 ordenações com os mesmos exemplos da estratégia mixed: class0\_first coloca todos da classe 0 primeiro, class1\_first inverte, shuffled embaralha aleatoriamente, e alternating intercala.

O resultado para seed=42 com 10-shot mostra diferenças pequenas: entre 88.6\% e 90.0\%. O delta máximo é de 1.4 pontos percentuais, o que é praticamente negligível.

É importante notar que este experimento é diferente do viés de ordem das classes no Slide 21: lá mudamos o texto do prompt, aqui mudamos a ordem dos exemplos fornecidos. O fato de o efeito ser pequeno aqui é reconfortante --- sugere que o LLM processa os exemplos de forma relativamente equitativa.
\end{frame}

% ============================================================
% SLIDE 32
% ============================================================
\begin{frame}{Slide 32 --- Validação Oráculo}
\small
A validação oráculo é nosso ``sanity check'' principal. A lógica é: se usarmos um classificador sintético perfeitamente consistente no lugar do LLM, os algoritmos deveriam recuperar exatamente o W verdadeiro.

Para o Problema D com o expert aniso\_x2, o Perceptron recupera W com similaridade cosseno de 0.9999 e fidelidade de 100\%. O NNLS fica em 0.9933 e 95.3\%. O LP também atinge 0.9999 e 100\%.

Isso é crucial para a interpretação: quando a fidelidade com o LLM é baixa, sabemos que o problema não está nos algoritmos --- está na inconsistência do LLM. Os algoritmos são capazes de resolver o problema quando o input é consistente. A inconsistência medida reflete propriedades reais do LLM, não limitações metodológicas.
\end{frame}

% ============================================================
% SLIDE 33
% ============================================================
\begin{frame}{Slide 33 --- R3 Metodologia}
\small
O experimento R3 testa se o LLM captura relações não-lineares entre features. A motivação é simular um kernel quadrático sem sair do mundo de métricas lineares: criamos uma terceira feature $x_3 = x_1 \cdot x_2$, que captura a interação entre as duas dimensões originais.

Os mesmos 150 pontos do Problema A são usados, agora com 3 coordenadas no prompt. Os 3 algoritmos aprendem um W com 3 componentes. A pergunta central é: o peso $w_3$ associado à interação é relevante? O LLM ``percebe'' a feature de interação?

Também comparamos as decisões do LLM com 2 features versus 3 features: se a concordância for baixa, significa que adicionar a feature extra muda o comportamento do modelo.
\end{frame}

% ============================================================
% SLIDE 34
% ============================================================
\begin{frame}{Slide 34 --- R3 Resultados}
\small
Os resultados são interessantes e ambíguos. O Perceptron atribui peso relevante a $w_3$ --- 1.07, comparável a $w_1$ = 2.31 e $w_2$ = 1.21. Isso sugeriria que a interação importa. Mas o NNLS atribui $w_3 \approx 0.03$, quase zero. E o LP distribui igualmente: 0.33 para cada.

A fidelidade R3 é boa: 86\% a 90\%, melhor que a fidelidade R2. Mas o dado mais revelador é a concordância cruzada: LLM com 2 features versus métrica com 3 features é de apenas 60\%. Isso mostra que adicionar $x_3$ ao prompt realmente muda o comportamento do LLM --- ele não ignora a feature extra.

A divergência entre algoritmos sobre $w_3$ é um resultado em si: mostra que a importância da feature de interação é mal determinada pelo conjunto de dados e depende da formulação do problema de otimização.
\end{frame}

% ============================================================
% SLIDE 35
% ============================================================
\begin{frame}{Slide 35 --- Baselines Configuração}
\small
Uma pergunta fundamental é: o LLM faz algo que um classificador trivial não faria com os mesmos dados? Para responder, treinamos 3 classificadores clássicos nos exatamente mesmos exemplos few-shot da Fase D.

O k-NN usa k ajustado ao número de exemplos, no máximo 5 e sempre ímpar, com distância euclidiana padrão. A Regressão Logística é linear com regularização L2. E o SVM usa kernel radial gaussiano com parâmetros padrão do sklearn.

A comparação é justa: mesmo treino, mesmo teste, mesmas métricas. A única diferença é o classificador. Se os baselines atingirem performance similar ao LLM, então o LLM não está fazendo nada especial --- está simplesmente classificando como qualquer outro algoritmo faria com aqueles dados.
\end{frame}

% ============================================================
% SLIDE 36
% ============================================================
\begin{frame}{Slide 36 --- Baselines Resultados}
\small
Os resultados são reveladores. Com exemplos easy e 10-shot, todos os classificadores --- incluindo o LLM --- ficam entre 88\% e 90\%. Os pontos são trivialmente separáveis e qualquer método funciona.

Mas com exemplos hard, a história é diferente. O k-NN despenca para 62\%, o SVM para 48\%, mas a Regressão Logística atinge 100\%! Isso faz sentido: a fronteira do expert é essencialmente linear, e a LR aprende exatamente essa fronteira. O k-NN e o SVM-RBF sofrem porque os poucos pontos na fronteira não fornecem vizinhança suficiente.

E o LLM? Fica em 68.5\% com hard --- melhor que k-NN e SVM, mas muito abaixo da LR. A conclusão é que o LLM não é uniformemente superior. Sua vantagem aparece em cenários de poucos exemplos ambíguos onde k-NN e SVM falham, mas a LR demonstra que algoritmos clássicos com a indução correta --- viés linear --- podem ser superiores.
\end{frame}

% ============================================================
% SLIDE 37
% ============================================================
\begin{frame}{Slide 37 --- Robustez entre Seeds}
\small
Olhando a robustez global entre as 3 seeds, a consistência média no Problema B é 81.1\% com desvio de 14.8\%, e no Problema C é 77.9\% com desvio de 16.8\%. A variabilidade é alta, especialmente em zero-shot.

A boa notícia é que os resultados few-shot são consideravelmente mais estáveis: o desvio cai significativamente com mais exemplos. Isso reforça que o LLM, quando orientado, é mais previsível.

A instabilidade de W entre seeds, já discutida anteriormente, é a principal fragilidade de H5. Porém, é importante notar que a instabilidade de W não necessariamente implica instabilidade das classificações --- métricas diferentes podem produzir classificações semelhantes em certas geometrias.
\end{frame}

% ============================================================
% SLIDE 38
% ============================================================
\begin{frame}{Slide 38 --- Testes Estatísticos}
\small
Para rigor estatístico, usamos Bootstrap com 10.000 reamostras para calcular intervalos de confiança de 95\%.

O teste pareado entre zero-shot e 10-shot mostra que a melhoria é estatisticamente significativa: delta de +0.322 no Problema B com intervalo [+0.257, +0.360] que exclui zero, e delta de +0.373 no Problema C com intervalo [+0.342, +0.400]. Os tamanhos de efeito, medidos por Cohen's d, são de 7.08 e 19.81 --- efeitos extremamente grandes.

Na Fase D, o teste Easy vs.\ Hard com 5-shot mostra delta de +0.231 com CI [+0.100, +0.340], confirmando estatisticamente que Easy supera Hard nos tamanhos menores de few-shot.

O teste da métrica estimada versus a euclidiana em zero-shot mostra diferença não significativa: CI inclui zero. Isso significa que, em zero-shot, a métrica aprendida não é significativamente melhor que a métrica euclidiana trivial.
\end{frame}

% ============================================================
% SLIDE 39
% ============================================================
\begin{frame}{Slide 39 --- Síntese das Hipóteses}
\small
Vamos à síntese final das hipóteses. H1 sobre consistência é parcialmente confirmada: Kappa vai de 0.14 em zero-shot para 0.93 com 10 exemplos. Consistência existe, mas depende de orientação.

H2 sobre few-shot é fortemente confirmada: o salto de 30 pontos percentuais com apenas 5 exemplos é um efeito robusto e reproduzível.

H3 sobre o LLM como aprendiz está confirmada: o LLM vai de 55\% a 96\% com exemplos do expert.

H4 sobre exemplos difíceis é refutada --- e consideramos que esta refutação é um dos achados mais valiosos do trabalho. LLMs se beneficiam mais de exemplos prototípicos que de exemplos fronteiriços.

H5 sobre estabilidade é parcial: W é instável entre seeds, mas classificações few-shot são estáveis.
\end{frame}

% ============================================================
% SLIDE 40
% ============================================================
\begin{frame}{Slide 40 --- Contribuições}
\small
As contribuições deste trabalho são quatro. Primeira: propomos um framework de otimização inversa aplicado a LLMs, com 3 algoritmos independentes validados por oráculo. Esta é uma abordagem nova que conecta a área de metric learning com a análise de decisões de modelos de linguagem.

Segunda: fornecemos evidência empírica sobre a consistência decisória de LLMs --- consistentes com orientação, inconsistentes sem.

Terceira: a refutação de H4 é um achado valioso para a comunidade de in-context learning. Mostramos que exemplos prototípicos são mais eficazes que exemplos fronteiriços para LLMs, o que contrasta com a intuição clássica de active learning.

Quarta: realizamos uma análise de sensibilidade abrangente que identifica confundidores --- viés de classe, semântica de features, template de prompt, ordem de exemplos --- que devem ser controlados em trabalhos futuros.
\end{frame}

% ============================================================
% SLIDE 41
% ============================================================
\begin{frame}{Slide 41 --- Limitações e Trabalhos Futuros}
\small
Quanto às limitações, a mais importante é que testamos apenas um modelo, o gpt-4o-mini. Não sabemos se os achados se generalizam para Claude, Gemini, ou GPT-4o. Segundo, usamos apenas dados sintéticos 2D --- dados reais e alta dimensão podem se comportar diferentemente. Terceiro, a métrica diagonal com $O(d)$ parâmetros é uma simplificação que pode não capturar interações entre features. A temperatura fixa em 0.0 também limita a generalização. E o LP frequentemente retornou infeasible, o que limita a utilidade desse algoritmo.

Para trabalhos futuros, os caminhos mais promissores são: testar múltiplos LLMs para separar propriedades de LLMs em geral de propriedades do gpt-4o-mini especificamente; usar dados reais como Iris e MNIST; explorar a Mahalanobis completa; investigar fronteiras não-lineares; e, talvez o mais impactante, aplicar esta abordagem para auditoria prática de decisões de IA em domínios críticos como saúde e justiça.
\end{frame}

% ============================================================
% SLIDE 42
% ============================================================
\begin{frame}{Slide 42 --- Agradecimentos}
\small
Gostaria de agradecer ao meu orientador pela paciência e pelas contribuições fundamentais ao direcionamento da pesquisa, à banca por aceitar avaliar este trabalho, e a todos que acompanharam esta jornada.

Fico à disposição para perguntas. Estou preparado para discutir qualquer aspecto da metodologia, dos resultados ou das implicações do trabalho.

Obrigado.
\end{frame}

\end{document}
